\documentclass[11pt,a4paper]{article}
\usepackage[utf8]{inputenc}
\usepackage[T1]{fontenc}
\usepackage{lmodern}
\usepackage[a4paper,top=1.5cm,bottom=1.5cm,left=2.4cm,right=2.4cm]{geometry}
\usepackage{microtype}
\usepackage[hidelinks]{hyperref}
\usepackage{parskip}
\pagestyle{empty}
\setlength{\parskip}{0.5em}

\begin{document}

\noindent\begin{flushright}
Olatunji Johnson\\
Department of Mathematics\\
University of Manchester\\
Manchester, UK\\
\texttt{olatunji.johnson@manchester.ac.uk}\\
24 June 2026
\end{flushright}

\noindent The Editors\\
\emph{Methods in Ecology and Evolution}\\
British Ecological Society

\medskip
\noindent Dear Editors,

I am pleased to submit my manuscript, \textbf{``Auditing spatial information leakage
in ecological prediction: quantifying the optimism in map accuracy with the
\texttt{spLeakage} R~package''}, for consideration as a Research Article / Application
in \emph{Methods in Ecology and Evolution}.

Ecologists increasingly turn field observations into wall-to-wall maps of species
occurrence, abundance, biomass and disease risk, and judge them by an accuracy
estimated through cross-validation. When data are spatially autocorrelated, ordinary
random cross-validation grades a model on near-interpolation rather than on the
extrapolation the map is actually used for, so reported accuracy is optimistically
inflated. A mature toolkit exists to \emph{prevent} this by building spatially
separated folds (including NNDM/kNNDM and \texttt{blockCV}, published in this
journal), but \textbf{no tool audits the validation an analyst has already used, or a
published map, and reports how inflated the accuracy is}. The field is moreover
divided over whether spatial cross-validation is even appropriate---a controversy
that has left practitioners without clear guidance.

My manuscript addresses both gaps with a single idea: spatial leakage and optimism
are well-defined only relative to a declared sampling \emph{design}, \emph{estimand}
and prediction \emph{target}. This reframing reconciles the controversy (each side is
correct in a different setting) and makes optimism computable. On this basis I
contribute (i) a \emph{signed} Spatial Leakage Index that audits any split and maps
where it leaks; (ii) a formal estimand for optimism under a Gaussian process, with a
closed form and a Wasserstein bound; (iii) a refit-free estimator of accuracy
inflation that flags when it is asked to extrapolate; and (iv) design-aware guidance,
a bias-corrected accuracy that can correct a published number, and optimal
validation-survey design. Everything is implemented in the open-source
\texttt{spLeakage} package, which \emph{audits rather than generates} folds and
interoperates with the existing ecosystem.

The work is a strong fit for \emph{MEE}: it builds on and complements methods this
journal has championed (NNDM, the area of applicability, \texttt{blockCV}), pairs a
methodological advance with documented open-source software, and speaks to emerging
reporting standards for distribution models. Key results: against exact
Gaussian-process algebra the index is a near-sufficient statistic for true optimism
($r=0.98$) and recovers the \emph{direction} of bias that unsigned matching statistics
cannot ($r=0.98$ vs $-0.09$); auditing 16 real datasets, reported accuracy is
optimistic in \emph{every} case (median 22\%); and two contrasting case studies---a
species distribution model and a malaria-prevalence map---show the audit correctly
attributing leakage to genuine autocorrelation in one and to a survey-design artefact
in the other.

The manuscript is original, is not under consideration elsewhere, and has not been
published previously; there are no conflicts of interest. The \texttt{spLeakage}
package and all analysis code are openly available at
\url{https://github.com/olatunjijohnson/spLeakage} (documentation at
\url{https://olatunjijohnson.github.io/spLeakage/}); all datasets are public, and code
will be archived with a DOI on acceptance. I have suggested several potential referees
spanning both the spatial cross-validation and design-based map-validation
perspectives the manuscript reconciles. Thank you for considering my submission.

\noindent Yours sincerely,\\[0.4em]
Olatunji Johnson

\end{document}
